\section{Preliminaries}\label{sec:preliminaries}

Throughout, $x$ tends to infinity.  Implied constants may depend on fixed
parameters and fixed seminorm bounds, but not on $x$ or on the particular
coefficient sequences under consideration.  We write
$e_q(z)=\exp(2\pi iz/q)$ and use $\tau$ for the divisor function.

\begin{definition}[Coefficient sequences]\label{def:coefficients}
A coefficient sequence is a function
$\alpha:\N\times(1,\infty)\to\R$ satisfying
\[
 |\alpha(n;x)|\ll \tau(n)^{O(1)}(\log x)^{O(1)}.
\]
It is \emph{located at scale} $N=N(x)$ if, for fixed $0<c<C$, it is
supported on $[cN,CN]$.  A sequence at scale $N$ has the
\emph{Siegel--Walfisz property} if, for all $A>0$, all $q,r\ge1$, and
all $(a,q)=1$,
\begin{equation}\label{eq:sw}
 \left|
 \sum_{\substack{n\equiv a\pmod q\\(n,r)=1}}\alpha(n;x)
 -\frac1{\varphi(q)}\sum_{(n,qr)=1}\alpha(n;x)
 \right|
 \ll_A \tau(qr)^{O(1)}\frac{N}{(\log x)^A}.
\end{equation}
It is \emph{smooth at scale} $N$ if
$\alpha(n;x)=\psi_x(n/N)$, where all $\psi_x$ are supported in one fixed
compact interval and every fixed derivative is bounded by a fixed power of
$\log x$.
\end{definition}

For $f=\alpha\star\beta$ put
\begin{equation}\label{eq:discrepancy}
 \Delta_f(a;d)=
 \sum_{\substack{x\le n\le2x\\n\equiv a\pmod d}}f(n;x)
 -\frac1{\varphi(d)}
 \sum_{\substack{x\le n\le2x\\(n,d)=1}}f(n;x).
\end{equation}
Thus \cref{eq:S52-conclusion} is the assertion
\[
 \sum_{\substack{d\in\DS\\d\ \mathrm{squarefree}}}
 |\Delta_f(a;d)|\ll_Ax(\log x)^{-A}.
\]

\subsection{Established analytic inputs}

We use two groups of previously established results.

\begin{enumerate}[label=\textup{(\roman*)}]
\item\label{input:trace} The sheaf-theoretic complete-sum and Fourier-transform estimates in
\cite[Theorem~6.17, Corollary~6.24, and Proposition~8.4]{PolymathEDZ2014}.
Only the second bound in Proposition~8.4 is used at the terminal stage.  Its
local rational phase is nondegenerate whenever the two denominator slopes are
distinct modulo every prime factor of the squarefree modulus.
\item\label{input:outer} The dispersion identity and initial Type-I reduction
in \cite[Sections~5 and 8A]{PolymathEDZ2014}, in the explicit form reorganized
for smooth moduli in \cite{Stadlmann2025}.  We use the identity before any
smoothness-based enlargement and prove the required enlargement for $\DS$.
\item\label{input:qvdc} The algebraic q-van der Corput identities from
\cite{Stadlmann2025}: Poisson completion, Cauchy--Schwarz, finite M\"obius
inversion, CRT substitutions, and Taylor expansion of fixed smooth profiles.
Every implication needed for the prescribed-factor family is stated and
proved below; the smooth-modulus conclusion itself is not used as an input.
\end{enumerate}

The Bombieri--Vinogradov theorem, the prime number theorem, and the standard
Maynard--Polymath multidimensional Selberg sieve are used in their customary
forms \cite{Maynard2015,PolymathSieve2014}.

\subsection{Parameter reserve}

Assume \cref{eq:S52-a,eq:S52-b,eq:S52-c}.  Set
\begin{align}
\mathfrak m=\min\{&1-8\omega-4\delta-2\gamma,
 \gamma-32\omega-10\delta,
 4\gamma-1-48\omega-16\delta,\notag\\
&\tfrac12-2\omega-\gamma,
 \gamma-\delta,
 \gamma-8\omega-4\delta,
 \tfrac14-16\omega-6\delta\}.
\label{eq:margin}
\end{align}
Every entry is positive.  We fix
\begin{equation}\label{eq:eta-choice}
 0<\eta<\eta_0:=
 \min\left\{\frac{\delta}{10^{100}},\frac{\mathfrak m}{1600}\right\},
 \qquad \eps_{\rm tr}=\frac{\eta}{1000}.
\end{equation}
The parameter $\eta$ pays structural losses in the q-van der Corput chain;
$\eps_{\rm tr}$ is reserved for conductor, divisor, and finite-seminorm
losses in the terminal trace estimate.  Keeping them separate prevents the
same power saving from being spent twice.

The first and third inequalities in the parameter region imply
\begin{equation}\label{eq:aux-range}
 12\omega+6\delta<\gamma<\frac12-2\omega.
\end{equation}
All later auxiliary range conditions follow from
\cref{eq:S52-a,eq:S52-b,eq:S52-c,eq:eta-choice}.
